\chapter{Implementation}

The implementation converts the system design into executable modules and generated artifacts. The system implements a reproducible workflow for dataset processing, spike encoding, PWL file generation and classifier integration. The implementation also defines how the generated PWL files are used by the Cadence/LTSpice LIF circuit design to evaluate neuromorphic performance.

\section{System Modules Organization}

The project implementation is organized into distinct processing modules. The logical separation of these modules ensures a clean pipeline from raw data to hardware simulation, as shown in Table \ref{tab:repo}.

\begin{table}[H]
\centering
\caption{System module organization}
\label{tab:repo}
\begin{tabular}{|p{0.35\textwidth}|p{0.5\textwidth}|}
\hline
\textbf{Module / Component} & \textbf{Purpose} \\
\hline
\textbf{Raw Data Storage} & Stores UNSW-NB15 network traffic records. \\
\hline
\textbf{Processed Features Data} & Stores PCA-reduced analog feature matrices. \\
\hline
\textbf{PWL Source Generator} & Stores generated PWL voltage source files for hardware simulation. \\
\hline
\textbf{Hardware Outputs} & Stores Cadence waveform exports and extracted feature outputs. \\
\hline
\textbf{Data Pipeline Module} & Implements preprocessing, normalization and PCA mapping. \\
\hline
\textbf{Spike Encoder Module} & Implements spike encoding and PWL file generation. \\
\hline
\textbf{Hardware Interface Module} & Provides routines for hardware waveform parsing and spike counting. \\
\hline
\textbf{Classification Module} & Provides classifier integration for extracted spike features. \\
\hline
\end{tabular}
\end{table}

\\section{Software Environment}

The core software implementation is developed in the Python 3.10 programming environment, leveraging industry-standard numerical and machine learning libraries to ensure high performance and reproducibility. 

Pandas is extensively utilized for large-scale data ingestion, cleaning, and matrix manipulation of the millions of network records. NumPy provides the highly optimized underlying $C$-based arrays necessary for rapid vector arithmetic during the spike encoding phase. The Scikit-learn (sklearn) library provides the robust, mathematically verified implementations for Min-Max scaling, label encoding, Principal Component Analysis (PCA) dimensionality reduction, and the final Support Vector Machine (SVM) classification models. Furthermore, specialized parsing routines are custom-built to interface natively with Cadence Virtuoso and LTSpice exported `.csv` raw simulation logs.

\section{Data Pipeline Module Implementation}

The Data Pipeline module rigorously implements dataset collection, deep preprocessing, and feature extraction. The module enforces a strict schema for the UNSW-NB15 dataset. All chunked dataset files matching the standard format are dynamically loaded from the raw dataset repository and concatenated into a singular, cohesive Pandas DataFrame.

\subsection{Derived Feature Engineering}
A critical step in the preprocessing pipeline is the derivation of temporal traffic intensity metrics. The implementation engineers a derived feature called \texttt{packet\_rate}. This is computed algorithmically using total source packets, destination packets, and the overall flow duration:
\begin{equation}
packet\_rate = \frac{Source\_Packets + Destination\_Packets}{Flow\_Duration + \epsilon}
\end{equation}
A minute epsilon constant ($\epsilon = 10^{-6}$) is mathematically injected when the flow duration is nominally zero, successfully averting catastrophic division-by-zero errors. This derived feature is highly sensitive to rapid network flooding attacks (e.g., DoS, DDoS).

\subsection{Data Cleansing and Normalization}
Network identifiers that would introduce severe model bias are systematically stripped from the feature matrix. The following specific columns are removed:
\begin{itemize}
\item Source and Destination IP addresses (prevents subnet bias).
\item Start and Last time temporal markers.
\item Binary label (as the multi-class attack category is strictly retained).
\end{itemize}

The target attack category column is retained and textually cleaned so that normal benign traffic is represented uniformly. Any remaining missing values (NaN) are filled via median imputation, and categorical fields (such as HTTP state or protocol) are strictly label encoded into integer identifiers. Finally, the entire continuous numerical matrix is mapped to the range $[0, 1]$ using a \texttt{MinMaxScaler}.

\section{PCA Mapping Implementation}

Because simulating analog circuits is computationally and temporally expensive, passing all 40+ raw network features into the hardware model is physically infeasible. Therefore, after initial normalization, PCA is applied to construct a dense, reduced representation. The algorithm is initialized as:
\begin{verbatim}
pca = PCA(n_components=8)
reduced_features = pca.fit_transform(normalized_data)
\end{verbatim}

The specific physical neuron count targeted in this circuit design is eight. The PCA linearly projects the high-dimensional network flow onto the eight vectors of maximal variance. Since PCA projection vectors inherently contain negative coordinate spaces and are completely unbounded, the Python script mathematically enforces a secondary, critical Min-Max normalization phase. This guarantees that all $N_{PCA}=8$ columns strictly reside within the $[0.0, 1.0]$ bounds expected by the analog-voltage spike encoder. 

The fully processed analog feature array is written to disk. Simultaneously, a smaller Hardware-in-the-Loop (HIL) subset containing a statistically representative sample across all nine attack categories is generated. This subset allows the ECE team to execute Cadence circuit simulations in a fraction of the time, avoiding the prohibitive computational cost of simulating transient analog dynamics for millions of dataset rows.

\section{Spike Encoder Module Implementation}

The Spike Encoder module is the highly mathematical core of the software-to-hardware translation layer. It actively converts static analog $[0, 1]$ values into dynamic temporal events. The script implements three distinct encoding algorithms:

\subsection{Rate and TTFS Encoding}
The \texttt{rate\_encoding(value)} function computes a proportional density of active spike ticks bounded by the simulation window. The \texttt{ttfs\_encoding(value)} (Time-To-First-Spike) function computes exactly one spike, where the latency of the spike is inversely proportional to the stimulus magnitude—meaning larger analog network anomalies trigger significantly faster physical hardware responses.

\subsection{Gaussian Population Coding}
The primary algorithm used for the final circuit interface is \texttt{population\_coding(value)}. For each of the 8 PCA features, $M=4$ overlapping Gaussian receptive fields are instantiated. The activation of the $i$-th population neuron is precisely computed as:
\begin{verbatim}
sigma = 1.0 / (M - 1)
activation = np.exp(-0.5 * ((value - mu) / sigma)**2)
\end{verbatim}
These Gaussian activations expand the representational capacity of the analog variables, ensuring high-fidelity signal mapping into the neuromorphic hardware space. 

\section{PWL Generation Implementation}

The function \texttt{translate\_to\_ltspice\_pwl()} is tasked with formatting the binary spike trains into raw time-voltage trajectory pairs that the Cadence Virtuoso analog simulation engine can parse flawlessly. 

To completely prevent standard SPICE simulator convergence failures (which occur when simulators encounter mathematically discontinuous, vertical voltage steps), the generation script algorithmically enforces a linear rise/fall slew rate equal to exactly 1\% of the defined timestep.

An excerpt of a generated Piecewise Linear (PWL) `.txt` file demonstrating this critical $1\text{ms}$ pulse width and $10\mu\text{s}$ engineered slew rate is shown below:
\begin{verbatim}
0.000000 0.0
0.000990 0.0
0.001000 1.0   <-- 1% Linear Rise
0.001990 1.0
0.002000 0.0   <-- 1% Linear Fall
0.002990 0.0
\end{verbatim}

\section{Hardware Interface Implementation}

The Hardware Interface module provides the mission-critical bridge between the analog Cadence circuit simulation outputs and the Python-based machine learning analysis. Cadence Virtuoso exports simulated transient waveforms as massive `.csv` files containing dense, raw X-Y column pairs representing microsecond-scale simulation steps.

\subsection{Waveform Parsing and Thresholding}
The Python parser dynamically identifies headers matching the required output traces (e.g., \texttt{/I8/vout Y} for the spike voltages and \texttt{/I8/vmembrane X} for the temporal axis). 

To cleanly extract binary spike events from the noisy analog output waveforms, the algorithm employs a strict threshold-crossing logic. It iterates over the continuous voltage array and registers a spike strictly when the output voltage $V_{out}$ positively crosses the $0.5\text{ V}$ hysteresis threshold. This rigorous edge-detection algorithm ensures that a single, wide analog hardware pulse is never erroneously counted multiple times by the machine learning algorithm.

\section{Classifier Integration Implementation}

The Classification Integration module serves as the final analytical endpoint for the entire project architecture. It receives the heavily processed arrays containing the physical hardware spike counts extracted from the Cadence waveform logs.

The final column of the array is safely parsed as the ground-truth IDS attack label, while the preceding columns represent the physical network features perfectly encoded as hardware spike counts. The implementation intelligently partitions the dataset into 80\% training and 20\% testing sets. 

A Support Vector Machine (SVM) utilizing a Radial Basis Function (RBF) kernel is deployed. The SVM algorithm is exceptionally well-suited for decoding the non-linear boundaries intrinsic to population-coded spike representations. The module concludes execution by computing and reporting comprehensive evaluation metrics, including overall diagnostic Accuracy, Weighted F1-scores, a detailed Confusion Matrix, and the crucial False Positive Rate, ultimately validating the security efficacy of the neuromorphic architecture.

\section{Cadence Virtuoso Usage Workflow}

The intended circuit-level usage workflow is highly structured to prevent data corruption between the software and hardware domains:
\begin{enumerate}
\item \textbf{Stimulus Selection}: The ECE operator navigates the generated \texttt{pwl\_sources/} directory and selects a specific network flow sample (e.g., \texttt{sample\_0\_cat\_6}).
\item \textbf{Circuit Integration}: The 32 individual PWL `.txt` files corresponding to the population-coded PCA features are attached to the independent input voltage sources within the Cadence Virtuoso LIF schematic.
\item \textbf{Transient Simulation}: The Analog Design Environment (ADE) transient simulation is executed for the total duration of the $10\text{ms}$ spike window, calculating the non-linear integration dynamics.
\item \textbf{Waveform Probing}: The operator explicitly probes the critical nodes: the input stimulus voltage, the analog membrane potential ($V_{memb}$), and the output spike comparator voltage ($V_{comp}$).
\item \textbf{Data Export}: The probed waveforms are exported as a high-resolution time-series `.csv` file.
\item \textbf{Python Integration}: The exported `.csv` is fed back into the Python \texttt{Hardware Interface Module} for automated spike extraction and final SVM evaluation.
\end{enumerate}

This modular workflow allows for both single-neuron isolated testing and fully time-multiplexed parallel simulation across multiple LIF cells.

\section{Implementation Traceability}

Traceability between the software and hardware domains is strictly maintained through systematic, algorithm-driven file generation and naming conventions. Because the analog simulation strips away the textual network data, preserving the provenance of each voltage signal is an absolute necessity.

The processed Python matrices perfectly preserve the original UNSW-NB15 attack category labels as integer identifiers. Consequently, the generated PWL sample directories explicitly encode both the sample index and the true class label in their naming scheme. This makes it possible to uniquely trace any physical circuit simulation output back to the exact network flow record that produced it. 

For example, a physical directory systematically designated as \texttt{sample\_0\_cat\_6/} unambiguously indicates to both the ISE and ECE teams that the contained PWL voltage sources represent dataset sample index 0, and that this sample corresponds to the encoded attack category 6. This eliminates any possibility of mislabeling the hardware outputs before they reach the SVM classifier.

\section{Quality Checks and System Verification}

To ensure complete robustness across the interdisciplinary boundary, the following automated checks are actively enforced to verify implementation correctness:
\begin{enumerate}
\item \textbf{Dimensionality Verification}: Confirm that the processed PCA features contain exactly eight neuron columns and one encoded label column, matching the physical circuit footprint.
\item \textbf{Bounds Checking}: Mathematically confirm that all secondary PCA feature values strictly lie within the $[0.0, 1.0]$ bounds required by the analog-voltage spike encoder.
\item \textbf{File Integrity}: Confirm that each generated PWL file contains valid, monotonic time-voltage pairs without any illegal vertical voltage steps that could crash the SPICE matrix solver.
\item \textbf{Digital Compliance}: Confirm that the discrete PWL voltage values toggle exclusively between $0.0\text{ V}$ and $1.0\text{ V}$.
\item \textbf{Population Completeness}: Confirm that each population-coded sample directory contains the expected 32 independent files ($8 \text{ features} \times 4 \text{ populations}$).
\item \textbf{Hysteresis Verification}: Confirm that the Python spike extraction module reliably counts discrete rising edges through the $0.5\text{ V}$ threshold, actively rejecting transient noise and repeated high samples.
\end{enumerate}

\section{Current Implementation Status}

The implemented project work successfully establishes the complete front-end software pipeline required for advanced neuromorphic IDS experimentation. The raw network dataset can be dynamically processed into dense analog PCA features, those features can be mathematically encoded into biological spike trains, and the resulting spike trains can be physically exported as SPICE-compatible PWL voltage files. Furthermore, the hardware interface and classifier modules precisely define the integration layer necessary for parsing analog circuit outputs and performing final IDS evaluation.

The present report rigorously documents the robust implementation foundation required for hardware validation. Detailed experimental results, physical circuit power measurements, and final analytical conclusions are intentionally outside the scope of this initial four-chapter design report.
